% Generated draft: requires mathematical and source review.
\documentclass[11pt]{article}
\usepackage[margin=25mm]{geometry}
\usepackage{fontspec}
\setmainfont{Latin Modern Roman}
\usepackage{amsmath,amssymb,mathtools,unicode-math}
\setmathfont{Latin Modern Math}
\usepackage{enumitem,needspace,xurl,hyperref}
\setlength{\parindent}{0pt}
\setlength{\parskip}{5pt}
\setlength{\emergencystretch}{3em}
\newcommand{\N}{\mathbb{N}}
\newcommand{\Z}{\mathbb{Z}}
\newcommand{\Q}{\mathbb{Q}}
\newcommand{\R}{\mathbb{R}}
\newcommand{\C}{\mathbb{C}}
\newcommand{\F}{\mathbb{F}}
\newcommand{\A}{\mathbb{A}}
\newcommand{\PP}{\mathbb{P}}
\newcommand{\E}{\mathbb{E}}
\newcommand{\eps}{\varepsilon}
\newcommand{\dd}{\,\mathrm{d}}
\begin{document}
\section*{UnsolvedMath ID: 30004547}
\textbf{Source problem code:} catalogue entry
\subsection{Definitions and mathematical statement}
Let \(M\) be a closed smooth real manifold, where closed means compact and without boundary. No connectedness assumption or dimension restriction beyond those required by the structures below is supplied. An almost-complex structure is a smooth bundle endomorphism \(J:TM\to TM\) satisfying \(J^2=-\operatorname{id}_{TM}\). It is integrable if \(M\) admits a complex-manifold atlas whose induced almost-complex structure is \(J\). A smooth real differential \(2\)-form \(\omega\) is nondegenerate if, at every \(p\in M\), the map \(T_pM\to T_p^*M\), \(v\mapsto\omega_p(v,\cdot)\), is an isomorphism. Either structure forces the dimension of each component to be even.

\textbf{Question.} Among closed manifolds admitting an almost-complex structure, determine additional necessary conditions on their topology for admitting an integrable almost-complex structure. More explicitly, which topological restrictions follow from the existence of an integrable \(J\), beyond those already implied by the existence of an almost-complex structure? The source asks broadly for such restrictions; it specifies neither a particular invariant nor a conjectured complete criterion. This concerns restrictions on the underlying manifold and does not assert that every almost-complex structure on an eligible manifold is integrable.

The original source formulation also asks the analogous question for a nondegenerate \(2\)-form \(\omega\), but does not define its ``integrability.'' Under the conventional interpretation that this means \(\dd\omega=0\), such a form is symplectic, and this branch asks: among closed manifolds admitting a nondegenerate \(2\)-form, which additional topological restrictions are necessary for the existence of a symplectic form? This interpretation is conditional and requires source confirmation. The two branches are alternatives; no compatibility between \(J\) and \(\omega\) is assumed.
\subsection{Short English statement}
Determine which additional restrictions on the topology of a closed manifold arise when it admits an integrable almost-complex structure. The original formulation also asks about nondegenerate \(2\)-forms; interpreting their integrability as closedness gives the analogous question for symplectic structures. The source leaves both the intended scope of the restrictions and the meaning of integrability for a \(2\)-form unspecified.
\subsection{Sources}
\begin{itemize}
\item Supplied local UnsolvedMath registry record, permanent ID \(30004547\), title: \emph{Topological Obstructions to Integrability of Geometric Structures}; Oberwolfach Reports Open Problems collection; exact source problem code \texttt{catalogue entry}. No individual proposer is supplied; the proposed year is recorded as 2020.
\item Supplied source citation: \emph{Mini-Workshop: Almost Complex Geometry (2021)}, Oberwolfach Reports, pp.\ 1657--1691, DOI \texttt{10.4171/owr/2020/33}. \url{https://doi.org/10.4171/owr/2020/33}. The literature-source list labels this reference \emph{Oberwolfach Report 33/2020}. The supplied citation's 2021 designation and the report's 2020 designation are retained without reconciliation.
\item Supplied literature reference: Albanese--Milivojevic, 2019. \url{https://arxiv.org/abs/1805.04751}. No title or more detailed bibliographic information is supplied.
\item Literature assessment dated \textbf{2026-08-22}: the supplied curation labels the problem open and states that broad topological restrictions imposed by integrability of \(J\) remain unresolved; it mentions constraints in special cases and for rational-homology spheres, and reports that no general classification of underlying manifolds is known. This is supplied dated curation, not an independently verified assessment. The background explicitly requests independent verification of time-sensitive status and cited claims.
\end{itemize}
\end{document}
